\chapter{作用量泛函与\personConley{}--\personZehnder{}指标}
\label{ch:action-conley-zehnder}

哈密顿\personFloer{}同调中有一种广为流传的看法：它是辛流形 $(M,\omega)$ 的自由环路
空间 $\mathcal L(M)$ 这一无限维空间上的“$(\infty/2)$ 维”同调理论。辛形式
$\omega$ 通过对环路取平均，自然地在 $\mathcal L(M)$ 上诱导一个微分形式
$\Omega$。定义域旋转诱导 $\mathcal L(M)$ 上的自然 $S^1$-作用，其生成向量场为
$\mathbb X(\gamma)=\dot\gamma$；这一作用又定义迭代积分
$\mathbb X\mathbin{\lrcorner}\Omega$。由于 $\Omega$ 具有 $S^1$-不变性，按照\personChen{}的迭代积分理论 (Chen73)、(GJP91)，形式
\stfhIndex{I026}{Zuoyongyixingshi}{作用 1-形式（action one-form）}
$\mathbb X\mathbin{\lrcorner}\Omega$ 是闭的。在辛几何中，这个闭 1-形式称为
\emph{作用 1-形式}，记作 $\alpha$。

当加入一个 1-周期哈密顿量 $H:S^1\times M\to\mathbb R$ 时，它的微分族
$\{dH_t\}_{t\in S^1}$ 在 $\mathcal L(M)$ 上诱导一个正合 1-形式，其原函数是平均函数
$\mathbb H:\mathcal L(M)\to\mathbb R$：
\[
  \mathbb H(\gamma):=\int_0^1 H(t,\gamma(t))\,dt.
\]
于是 $\alpha+d\mathbb H$ 仍是闭 1-形式（在 (Chen73) 的意义下）。$\alpha$ 是
$S^1$-不变的，而除非 $\mathbb H$ 是自治哈密顿量，$d\mathbb H$ 并不具有这一
不变性。\footnote{译者注：源书此句主语作黑板体 $\mathbb H$（即上式的平均
泛函），并称其为“自治哈密顿量”；$\mathbb H$ 并非哈密顿量，其微分
$d\mathbb H$ 的 $S^1$-不变性当且仅当 $H$ 自治。译文照录源书字形，特此说明。}
\stfhIndex{I137}{Nuoweikefu}{诺维科夫（Novikov）}
\personNovikov{} (No81, No82) 在有限维情形发展了这类闭 1-形式的\personMorse{}理论。从这个意义
说，可以把\personFloer{}同调看作 1-形式 $\alpha+d\mathbb H$ 的\personNovikov{}同调。经典作用量
泛函在 $\mathbb C^n$ 上的临界点理论，例如 (Bn82, BnR79) 中的发展，被\personNovikov{}--\personFloer{}理论的椭圆方法取代；在过去三十年里，这一直是推动辛拓扑和哈密顿
动力学发展的强大动力。

\stfhIndex{I091}{Ziyouhuanlukongjian}{自由环路空间（free loop space）}
\section{\texorpdfstring{自由环路空间及其 $S^1$-作用}{自由环路空间及其 S¹-作用}}
\label{sec:free-loop-space-circle-action}

设 $M$ 是一般光滑流形，不必是辛流形。记
$\mathcal L(M):=\operatorname{Map}(S^1,M)$ 为\emph{自由环路空间}，即光滑映射
\[
  \gamma:S^1=\mathbb R/\mathbb Z\longrightarrow M
\]
的集合。我们强调，环路带有标记点 $0\in\mathbb R/\mathbb Z$；我们也常以单位区间
$[0,1]$ 参数化环路，并附加周期边界条件 $\gamma(0)=\gamma(1)$。
$\mathcal L(M)$ 有一个由可缩环路组成的特定连通分支，记作 $\mathcal L_0(M)$。
其\termUniversalCoveringSpace{}，记作 $\widetilde{\mathcal L}_0(M)$，可以写成丛
$\pi:\widetilde{\mathcal L}_0(M)\to\mathcal L_0(M)$，在
$\gamma\in\mathcal L_0(M)$ 上的纤维为
\[
 \{[\gamma,w]\mid w:D^2\to M,\ \partial w:=w|_{\partial D^2}=\gamma\},
\]
其中 $[\gamma,w]$ 表示满足 $\partial w=\gamma$ 的圆盘
$w:D^2\to M$ 相对于边界的同伦类。这里通过乘以 $2\pi$ 将 $\partial D^2$ 与
$S^1$ 认同。称这样的 $w$ 为 $\gamma$ 的\emph{填充圆盘}。

由纤维化 $\Omega_0(M)\to\mathcal L(M)\to M$ 可得
\[
 \pi_1(\mathcal L(M))\cong\pi_1(\Omega_0(M))\times\pi_1(M),
\]
其中 $\Omega_0(M)$ 是基环路空间，比如以常环路
$\widetilde y_0\in\mathcal L_0(M)$ 为基点，且 $y_0\in M$。回忆
$\pi_1(\Omega_0(M))=\pi_1(\Omega_0(M),\widetilde y_0)\cong\pi_2(M)$，便得到
\footnote{译者注：原书在这段纤维化中仍以 $\mathcal L_0(M)$ 表示基环路空间，
与前文同一符号所表示的可缩自由环路分支冲突；译文以常用记号 $\Omega_0(M)$ 区分
基环路空间，而式 (18.1.1) 仍照录原书的群同构。}
\footnote{译者注：在一般光滑流形这一假设下，评价纤维化通常只给出
$\pi_2(M)$ 被 $\pi_1(M)$ 的单值性作用扭曲的半直积，而不一定是式 (18.1.1) 所写的
直积；仅以填充圆盘为数据的覆盖也不一定是万有覆盖。原书未补充使这些断言成立的
额外假设，故正文照录并在此提示。}
\begin{equation}
 \pi_1(\mathcal L_0(M),\widetilde y_0)
 \cong\pi_2(M,y_0)\times\pi_1(M,y_0).
 \label{eq:18-free-loop-fundamental-group}
\end{equation}

与元素 $(A,a)\in\pi_2(M)\times\pi_1(M)$ 相伴的
$\widetilde{\mathcal L}_0(M)$ 的\termDeckTransformation{}可如下实现。以满足
$v(N)=y_0$ 的球面 $v:S^2\to M$ 表示 $A$，并取
$[\gamma,w]\in\widetilde{\mathcal L}_0(M)$。把区间 $[0,1]$ 看作
$D^2$ 中的子集 $[0,1]\times i\cdot\{0\}$，并以 $0$ 表示 $D^2$ 的原点。
任选道路 $\ell_1:[0,1]\to M$，使得
$\ell_1(0)=y_0$、$\ell_1(1)=w(0)$，把 $w(0)$ 与 $y_0$ 连接起来。
\footnote{译者注：源书印 $\ell_1(0)=w(0)$、$\ell_1(1)=0$，端点互换且基点误写为
$0$；译文补正为 $\ell_1(0)=y_0$、$\ell_1(1)=w(0)$，使拼接
$\ell_1\#w|_{[0,1]}\#\overline{\ell}_2$ 成为以 $y_0$ 为基点的环路。}
再选道路
$\ell_2:[0,1]\to M$，满足
\[
 \ell_2(0)=y_0,\qquad \ell_2(1)=\gamma(1),
\]
并考虑连接
\[
 \ell_1\#w|_{[0,1]}\#\overline{\ell}_2,
\]
它给出一个以 $y_0$ 为基点的环路。要求这个环路的同伦类就是给定的
$a\in\pi_1(M)$。特别地，$A\in\pi_2(M)$ 在 $[\gamma,w]$ 上的作用由“粘上一个
球面”的操作实现：
\begin{equation}
 (\gamma,w)\longmapsto(\gamma,w\#v),
 \label{eq:18-glue-sphere-action}
\end{equation}
其中模去上述 $\pi_1(M)$ 的作用，而 $v:S^2\to M$ 是表示 $A$ 的某个（因而任意一个）
球面。

$\mathcal L(M)$ 上有由定义域平移诱导的自然圆周作用
\begin{equation}
 \gamma\longmapsto\gamma\circ R_\varphi=\gamma(\,\cdot+\varphi),
 \label{eq:18-loop-rotation}
\end{equation}
其中 $R_\varphi:S^1\to S^1$ 定义为
\[
 R_\varphi(t)=t+\varphi,\qquad\varphi\in S^1.
\]
这一作用的无穷小生成元是 $\mathcal L(M)$ 上的向量场 $\mathbb X$：
\begin{equation}
 \mathbb X(\gamma)=\dot\gamma.
 \label{eq:18-circle-generator}
\end{equation}
该 $S^1$-作用的不动点集是常环路的集合
\[
 M\lhook\joinrel\longrightarrow\mathcal L_0(M).
\]
这一作用还提升为二元组 $(\gamma,w)$ 集合上的作用
\begin{equation}
 (\gamma,w)\longmapsto(\gamma\circ R_\varphi,w\circ R_\varphi),
 \label{eq:18-lifted-circle-action}
\end{equation}
其中再次用 $R_\varphi$ 表示复乘法所诱导的映射
\[
 R_\varphi:z\in D^2\subset\mathbb C\longmapsto e^{2\pi i\varphi}z.
\]

\begin{lemma}
\label{lem:18-fixed-points-lifted-circle}
$\widetilde{\mathcal L}_0(M)$ 上诱导的 $S^1$-作用，其不动点集由
$[\gamma_x,w_x]$ 组成，其中 $\gamma_x$ 和 $w_x$ 分别是像为 $x$ 的常环路和常圆盘。
\footnote{译者注：按前文的万有覆盖模型，常环路也可配以代表非零球面类的填充圆盘；
同样利用 $R_\varphi\simeq\mathrm{id}$，这类加帽常环路的同伦类似乎也为不动点。
原书本引理只列常圆盘，与第 18.2 节稍后
$\operatorname{Crit}\mathcal A_0=\bigsqcup_{g\in\Gamma}g\cdot M$ 的描述亦有张力；
因缺少进一步约定，译文照录本引理并在此注明。}
\end{lemma}

\begin{proof}
定义域旋转 $R_\varphi$ 在常环路 $\gamma\equiv x\in M$ 上作用平凡。另一方面，
它在二元组 $(\gamma_x,w_x)$ 的同伦类上也作用平凡，因为 $R_\varphi$ 同伦于恒等映射，
所以 $(\gamma_x,w_x)$ 与
$(\gamma_x\circ R_\varphi,w_x\circ R_\varphi)$ 彼此同伦。
\end{proof}

\section{辛流形的自由环路空间}
\label{sec:free-loop-space-symplectic-manifold}

现在把上一节的讨论专门化到辛流形 $(M,\omega)$ 的情形。

\subsection{\texorpdfstring{$S^1$-作用及其矩映射}{S¹-作用及其矩映射}}
\label{subsec:circle-action-moment-map}

此时 $\mathcal L(M)$ 带有如下定义的典范（弱）辛形式 $\Omega$：
\begin{equation}
 \Omega(\xi_1,\xi_2):=\int_0^1\omega(\xi_1(t),\xi_2(t))\,dt.
 \label{eq:18-loop-space-symplectic-form}
\end{equation}
我们概述这个 2-形式的一些基本众所周知性质，而不试图给出完全严格的证明。关于一般
环路空间上的微分形式，可参见 (Chen73)、(GJP91) 的严谨论述。

\begin{proposition}
\label{prop:18-loop-space-symplectic-properties}
\begin{enumerate}[label=(\arabic*)]
\item 形式 $\Omega$ 是弱非退化的。换言之，对任意 $\gamma\in\mathcal L(M)$，
若对所有 $\eta\in T_\gamma\mathcal L(M)$ 都有 $\Omega(\xi,\eta)=0$，则且仅则
$\xi=0$。
\item $\Omega$ 是闭的。
\item 式~\eqref{eq:18-loop-rotation} 的 $S^1$-作用是辛的，即保持
$\Omega$；也就是说 $\mathcal L_{\mathbb X}\Omega=0$。
\footnote{译者注：原书此项引用式 (18.1.5)，但该式是覆盖空间上二元组
$(\gamma,w)$ 的提升作用；本命题讨论 $\mathcal L(M)$ 上的 $\Omega$，相应的定义域
旋转是式 (18.1.3)，译文据此改正引用。}
\end{enumerate}
\end{proposition}

$\Omega$ 的（弱）非退化性来自 $\omega$ 的非退化性；$\Omega$ 的闭性来自
$\omega$ 的闭性以及 \emph{$S^1$ 没有边界这一事实}。有关计算留给读者，也可参见 (Wn78)
对这些陈述的验证。

于是，通过投影 $\widetilde{\mathcal L}_0(M)\to\mathcal L_0(M)$ 的拉回，形式
$\Omega$ 在覆盖空间 $\widetilde{\mathcal L}_0(M)$ 上诱导一个辛形式，记为
$\widetilde\Omega$。

\stfhIndex{I050}{Dianfanbozhumian}{典范薄柱面（canonical thin cylinder）}
为说明上面陈述 (3) 的含义，先回忆两个邻近环路之间的\emph{典范薄柱面}这一有用
概念。记 $\exp$ 为度量
\[
 g:=\omega(\,\cdot,J_{\mathrm{ref}}\,\cdot)
\]
的指数映射，其中 $J_{\mathrm{ref}}$ 是任意给定的固定相容殆复结构。记
$\iota(g)$ 为 $g$ 的\termInjectivityRadius{}。只要给定的两点满足
$d(x,y)<\iota(g)$，就可以唯一地写成
\[
 y=\exp_x(E(x,y)),\qquad E(x,y):=(\exp_x)^{-1}(y).
\]
因此，如果环路 $\gamma,\gamma':S^1\to M$ 之间的 $C^0$ 距离
$d_{C^0}(\gamma,\gamma')$ 小于 $\iota(g)$，就可以定义典范映射
\[
 u^{\mathrm{can}}_{\gamma\gamma'}:[0,1]\times S^1\longrightarrow M
\]
为
\begin{equation}
 u^{\mathrm{can}}_{\gamma\gamma'}(s,t)
 =\exp_{\gamma(t)}\bigl(sE(\gamma,\gamma')(t)\bigr),
 \label{eq:18-canonical-thin-cylinder}
\end{equation}
其中 $E(\gamma,\gamma')(t):=(\exp_{\gamma(t)})^{-1}(\gamma'(t))$。

\begin{lemma}
\label{lem:18-action-one-form-locally-exact}
设 $\mathbb X$ 是式~\eqref{eq:18-circle-generator} 中的向量场。形式
$\mathbb X\mathbin{\lrcorner}\Omega$ 是 $\mathcal L(M)$ 上的闭 1-形式；具体说，
在任意给定环路 $\gamma_0$ 的一个 $C^1$ 邻域上，都存在函数
$\mathcal A=\mathcal A_0$，满足
\begin{equation}
 d\mathcal A=\mathbb X\mathbin{\lrcorner}\Omega.
 \label{eq:18-local-action-primitive}
\end{equation}
\end{lemma}

\begin{proof}
对任意与给定 $\gamma_0$ 在 $C^\infty$ 意义下充分接近的环路 $\gamma$，考虑
\[
 u^{\mathrm{can}}_{\gamma_0\gamma}:s\in[0,1]
 \longmapsto\exp_{\gamma_0}\bigl(sE(\gamma_0,\gamma)\bigr),
\]
它定义了端点固定的道路的一个特定同伦类
$[u_{\gamma_0\gamma}]$，且 $u(0)=\gamma_0$、$u(1)=\gamma$。相伴的参数化柱面
\[
 [0,1]\times S^1\longrightarrow M,\qquad
 (s,t)\longmapsto\exp_{\gamma_0(t)}\bigl(sE(\gamma_0(t),\gamma(t))\bigr)
\]
也记作 $u^{\mathrm{can}}_{\gamma_0\gamma}$。从它的显式表达式可知，在
$\gamma_0\in\mathcal L(M)$ 的一个 $C^1$-小邻域中，可以用公式
\begin{equation}
 \mathcal A(\gamma;\gamma_0)=0-\int
   (u^{\mathrm{can}}_{\gamma_0\gamma})^*\omega
 \label{eq:18-local-action-function}
\end{equation}
定义函数 $\mathcal A$。这里的“$0$”应理解为可任意选取的值
$\mathcal A(\gamma_0;\gamma_0)$。

下面验证~\eqref{eq:18-local-action-primitive}。设 $\gamma$ 位于给定的
$C^1$-小邻域中，$\xi\in T_\gamma\mathcal L(M)$。用道路芽
$\{\gamma_s\}_{-\epsilon<s<\epsilon}$ 表示 $\xi$，并对
$-\epsilon<s<\epsilon$ 考虑二元组 $(\gamma_s,u_s)$，其中
\[
 u_s:=u^{\mathrm{can}}_{\gamma_0\gamma}\#\{\gamma_s\}_{[0,s]}.
\]
计算得
\begin{align}
 d\mathcal A(\gamma)(\xi)
 &=\left.\frac d{ds}\right|_{s=0}\left(-\int u_s^*\omega\right)\notag\\
 &=\left.\frac d{ds}\right|_{s=0}\left(
   -\int (u^{\mathrm{can}}_{\gamma_0\gamma})^*\omega
   -\int_0^s\int_0^1
   \omega\left(\frac{\partial\gamma_u(t)}{\partial u},
                \frac{\partial\gamma_u(t)}{\partial t}\right)dt\,du
   \right)\notag\\
 &=\int_0^1-\omega(\xi(t),\dot\gamma(t))\,dt
  =(\mathbb X(\gamma)\mathbin{\lrcorner}\Omega)(\xi).
 \label{eq:18-action-one-form-computation}
\end{align}
这与~\eqref{eq:18-action-one-form-computation} 一起证明了
\eqref{eq:18-local-action-primitive}，从而证明引理。
\end{proof}

\begin{remark}
\label{rem:18-loop-space-de-rham-class}
从环路空间德拉姆理论 (Chen73)、(GJP91) 的观点看，$M$ 上的辛形式 $\omega$
通过闭 1-形式 $\mathbb X\mathbin{\lrcorner}\Omega$ 诱导一个典范的 1 次上同调类。
\stfhIndex{I115}{Diedaijifen}{迭代积分（iterated integral）}
用 (GJP91) 第 344 页的记号，该类由\emph{迭代积分} $\widetilde P_1(\omega)$ 得到。这个
\stfhIndex{I196}{Ruozhenghe}{弱正合（weakly exact）}
1-形式一般\emph{不是正合}的；它正合，恰好就是辛形式 $\omega$ 所谓的\emph{弱正合性}。
\end{remark}

若把这个闭 1-形式限制到 $\mathcal L_0(M)$，再考虑它在万有覆盖空间
$\widetilde{\mathcal L}_0(M)$ 上的提升，则公式~\eqref{eq:18-local-action-function}
可由二元组 $(\gamma,w)$ 的函数推广为\emph{全局}提升，仍记作
$\mathcal A=\mathcal A_0$：
\[
 \mathcal A_0(\gamma,w)=-\int w^*\omega.
\]
这里把 $w$ 看作从常道路 $w(0)$ 到 $\gamma=\partial w$ 的道路。积分前取负号，是为了
与作者关于谱不变量的一系列论文（例如 (Oh06a)）所用约定一致。称
$\mathcal A_0$ 为\emph{未扰动作用量泛函}。它在
$\widetilde{\mathcal L}_0(M)$ 上满足
\begin{equation}
 d\mathcal A_0=\mathbb X\mathbin{\lrcorner}\Omega.
 \label{eq:18-global-action-primitive}
\end{equation}
换言之，$\widetilde{\mathcal L}_0(M)$ 上的 $S^1$-作用是哈密顿的，与之相伴的矩映射
正是函数 $\mathcal A_0:\widetilde{\mathcal L}_0(M)\to\mathbb R$。更详细的讨论见
(Wn78)。

\subsection{\personNovikov{}覆盖}
\label{subsec:novikov-covering}

万有覆盖空间虽提供了使作用 1-形式的提升成为正合形式的模型，却太大而无法在其上
开展有意义的临界点理论：相伴函数的对称群
$\pi_1(\mathcal L_0(M))\cong\pi_2(M)\times\pi_1(M)$（参见
式~\eqref{eq:18-free-loop-fundamental-group} 后的译者注）可能高度\emph{非交换}，除非 $M$
单连通。因此要寻找一个更小的覆盖空间，使其覆盖变换群较小且为交换群。\personNovikov{}
(No81) 最先通过考虑一个使闭 1-形式的拉回成为正合形式的\emph{循环}覆盖，发展了紧流形上
闭 1-形式的\personMorse{}理论，并用 $M$ 的拓扑数据估计闭 1-形式的零点数。

\stfhIndex{I138}{Nuoweikefufugai}{诺维科夫覆盖（Novikov covering）}
依照 (Fl89b)、(HS95)，现在引入 $\mathcal L_0(M)$ 的\emph{\termNovikovCovering{}}概念。
它是 $\mathcal L_0(M)$ 的一个交换覆盖空间，对应于\personNovikov{} \personMorse{}理论中的循环覆盖。
沿用\personSeidel{} (Se97) 的术语，作如下定义。

\begin{definition}
\label{def:18-gamma-equivalence}
设 $(\gamma,w)$ 是由 $\gamma\in\mathcal L_0(M)$ 和以 $\gamma$ 为边界的圆盘 $w$
\stfhIndex{I007}{00-gamma-equivalence}{$\Gamma$-等价关系（$\Gamma$-equivalence）}
组成的二元组。称 $(\gamma,w)$ 与 $(\gamma,w')$ \emph{$\Gamma$-等价}，当且仅当
\[
 \omega([w'\#\overline w])=0,\qquad
 c_1([w'\#\overline w])=0,
\]
其中 $\overline w$ 表示在定义域上取相反定向的映射，$w'\#\overline w$ 是显然的
粘合球面。这里 $\Gamma$ 表示群
\[
 \Gamma=\frac{\pi_2(M)}
 {\ker(\omega|_{\pi_2(M)})\cap\ker(c_1|_{\pi_2(M)})}.
\]
记
\[
 \Gamma_\omega:=\omega(\Gamma)=\omega(\pi_2(M))\subset\mathbb R,
\]
\stfhIndex{I151}{Zhouqiqun}{周期群（period group）}
并称它为 $(M,\omega)$ 的（球面）\emph{\termPeriodGroup{}}。
\end{definition}

\begin{definition}
\label{def:18-rational-symplectic-manifold}
\stfhSubIndex{I186}{Xinliuxing}{辛流形（symplectic manifold）}{Youli}{有理（rational）}
若 $\Gamma_\omega\subset\mathbb R$ 是离散子群，则称 $(M,\omega)$ 为\emph{有理的}；
\stfhSubIndex{I185}{Xinliuxing}{辛流形（symplectic manifold）}{Wuli}{无理（irrational）}
否则称其为\emph{无理的}。
\end{definition}

\begin{example}
\label{ex:18-product-spheres-rationality}
带有乘积辛形式 $\omega_1\oplus\omega_2$ 的乘积
$S^2(r_1)\times S^2(r_2)$ 是有理的，当且仅当比值 $r_2^2/r_1^2$ 是有理数。
\end{example}

\begin{remark}
\label{rem:18-irrational-period-group}
注意，对无理的 $(M,\omega)$，周期群是 $\mathbb R$ 的可数稠密子集。一般认为，无理
辛流形上的哈密顿流，其动力学行为会比有理辛流形上的情形复杂得多。周期群
$\Gamma_\omega$ 是两者这种不同动力学行为的最简单指标。
\end{remark}

从现在起，专用 $[\gamma,w]$ 表示 $(\gamma,w)$ 的 $\Gamma$-等价类，并以
$\widetilde{\mathcal L}_0(M)$ 表示这些等价类的空间。记典范投影为
$\pi:\widetilde{\mathcal L}_0(M)\to\mathcal L_0(M)$，称
$\widetilde{\mathcal L}_0(M)$ 为 $\mathcal L_0(M)$ 的\termGammaCoveringSpace{}。
对 $A\in\pi_2(M)$，用 $A$ 或 $q^A$ 表示它在投影
$\pi_2(M)\to\Gamma$ 下的像。与 $A$ 相伴有两个自然不变量：\emph{赋值} $\nu(A)$，
\begin{equation}
 \nu:\Gamma\longrightarrow\mathbb R,\qquad \nu(A)=\omega(A),
 \label{eq:18-gamma-valuation}
\end{equation}
以及\emph{次数} $d(A)$，
\begin{equation}
 d:\Gamma\longrightarrow\mathbb Z,\qquad d(A)=2c_1(A).
 \label{eq:18-gamma-degree}
\end{equation}
一般而言，这两个不变量相互独立（见 (Gom98)），所以 $q^A$ 是依赖于两个变量的
形式参数。从这个意义说，也可以写成
\[
 q^A=T^{\omega(A)}e^{c_1(A)},
\]
其中 $T$ 与 $e$ 是两个不同的形式参数。规定 $T$ 的次数为 $0$，$e$ 的次数为 $2$。
\footnote{译者注：原书把 $e$ 的指数印作 $c_1(A)/2$，同时规定
$d(A)=2c_1(A)$ 且 $e$ 的次数为 $2$，两者给出的 $q^A$ 次数不一致；下一章又明确使用
$q^A=T^{\omega(A)}e^{c_1(A)}$。译文按次数定义及下一章记号补正。}

上面定义的（未扰动）作用量泛函 $\mathcal A_0$ 显然可下降到 $\Gamma$-覆盖空间，
公式仍为
\[
 \mathcal A_0([\gamma,w])=-\int w^*\omega.
\]
这个泛函在 $\widetilde{\mathcal L}_0(M)$ 上给出一个自然递增滤过：对每个
$\lambda\in\mathbb R$，定义
\[
 \widetilde{\mathcal L}_0^\lambda(M):=
 \{[z,w]\in\widetilde{\mathcal L}_0(M)\mid
   \mathcal A_0([z,w])\leq\lambda\}.
\]
注意，当 $\lambda\leq\lambda'$ 时，
\[
 \widetilde{\mathcal L}_0^\lambda(M)
 \subset\widetilde{\mathcal L}_0^{\lambda'}(M).
\]
由~\eqref{eq:18-global-action-primitive}，$\mathcal A_0$ 的临界集
$\operatorname{Crit}\mathcal A_0$ 是若干份 $M$ 的不交并：
\[
 \operatorname{Crit}\mathcal A_0=\bigsqcup_{g\in\Gamma}g\cdot M,
\]
其中 $M\hookrightarrow\mathcal L_0(M)$，$x\mapsto[x,\widetilde x]$ 是典范嵌入，
$[x,\widetilde x]$ 是常环路 $x$ 与常圆盘 $\widetilde x$ 组成的二元组。我们有自然交换图
\[
 \begin{array}{ccc}
  \operatorname{Crit}\mathcal A_0&\longrightarrow&
    \widetilde{\mathcal L}_0(M)\\
  \big\downarrow\scriptstyle\pi&&\big\downarrow\scriptstyle\pi\\
  M&\longrightarrow&\mathcal L_0(M),
 \end{array}
\]
其公共纤维同构于 $\Gamma$。

\subsection{第二变分}
\label{subsec:unperturbed-action-second-variation}

现在计算 $\mathcal A_0$ 在每个临界点 $[x,\widetilde x\#A]$ 处的第二变分
$d^2\mathcal A_0$。利用覆盖投影的切映射所诱导的典范认同
\begin{equation}
 T_{[x,\widetilde x]}\widetilde{\mathcal L}(M)
 \cong T_x\mathcal L(M),
 \label{eq:18-cover-tangent-identification}
\end{equation}
并把 $d^2\mathcal A_0([x,\widetilde x\#A])$ 看作定义在
$T_x\mathcal L(M)$ 上的二次型。

\begin{proposition}
\label{prop:18-unperturbed-action-bott-morse}
在每个 $[x,\widetilde x\#A]\in\operatorname{Crit}\mathcal A_0$ 处，Hessian
$d^2\mathcal A_0$ 在
\[
 T_{[x,\widetilde x\#A]}\widetilde{\mathcal L}_0(M)
 \cong T_x\mathcal L_0(M)
\]
上定义一个双线性型；它在 $\mathcal L_0(M)$ 中
$\operatorname{Crit}\mathcal A_0$ 的法方向上（弱）非退化。
\end{proposition}

\begin{proof}
为简化记号，记
\[
 [x,\widetilde x]=\widetilde x,\qquad
 [x,\widetilde x\#A]=\widetilde x\otimes q^A.
\]
由于作用 $\#A$ 是覆盖空间
$\widetilde{\mathcal L}_0(M)\to\mathcal L_0(M)$ 的覆盖变换，只需考虑点
$[x,\widetilde x]$。设
$V:(-\epsilon,\epsilon)\times(-\epsilon,\epsilon)\to\mathcal L(M)$ 是二参数环路族，
满足
\[
 \begin{cases}
 V(0,0,t)\equiv x,\\
 \left.\dfrac{\partial V}{\partial s}\right|_{s=u=0}(t)=\xi_1(t),\quad
 \left.\dfrac{\partial V}{\partial u}\right|_{s=u=0}(t)=\xi_2(t).
 \end{cases}
\]
以 $\gamma_{u,s}:S^1\to M$ 表示环路
$\gamma_{u,s}(t):=V(u,s,t)$，并令 $w_{u,s}:D^2\to M$ 为满足
$w_{u,s}|_{\partial D^2}=\gamma_{u,s}$ 的“小”圆盘。这样的圆盘在同伦意义下唯一。

由第二变分的定义，
\[
 d^2\mathcal A_0([x,\widetilde x])(\xi_1,\xi_2)
 =\left.\frac{\partial^2}{\partial u\,\partial s}\right|_{s=u=0}
  \mathcal A_0([\gamma_{u,s},w_{u,s}]).
\]
这里通过认同~\eqref{eq:18-cover-tangent-identification}，把 $\xi_i$ 看成
$T_x\mathcal L(M)$ 中的切向量。由~\eqref{eq:18-global-action-primitive}，
\[
 \frac\partial{\partial s}\mathcal A_0([\gamma_{u,s},w_{u,s}])
 =-\int\omega\left(\frac{\partial V}{\partial t},
                     \frac{\partial V}{\partial s}\right)dt.
\]
因此，使用任意无挠联络 $\nabla$，得到
\begin{align*}
 &\left.\frac{\partial^2}{\partial u\,\partial s}\right|_{s=u=0}
 \mathcal A_0([\gamma_{u,s},w_{u,s}])\\
 &=-\int_0^1(\nabla_u\omega)(x)
   \left(\left.\frac{\partial V}{\partial t}\right|_{u=s=0},
         \left.\frac{\partial V}{\partial s}\right|_{u=s=0}\right)dt\\
 &\quad-\int_0^1\omega\left(
   \left.\nabla_u\frac{\partial V}{\partial t}\right|_{u=s=0},
   \left.\frac{\partial V}{\partial s}\right|_{u=s=0}\right)dt
 -\int_0^1\omega\left(
   \left.\frac{\partial V}{\partial t}\right|_{u=s=0},
   \left.\nabla_u\frac{\partial V}{\partial s}\right|_{u=s=0}\right)dt.
\end{align*}
因为 $V(0,0,t)\equiv x$，故
$\left.\partial V/\partial t\right|_{u=s=0}\equiv0$，第一项和第三项均消失。对第二项，
由 $\nabla$ 无挠可得
\[
 \nabla_u\frac{\partial V}{\partial t}
 =\nabla_t\frac{\partial V}{\partial u}.
\]
从而
\[
 \left.\nabla_u\frac{\partial V}{\partial t}\right|_{u=s=0}
 =\left.\nabla_t\frac{\partial V}{\partial u}\right|_{u=s=0}
 =\left.\nabla_t\xi_2\right|_{u=s=0}
 =\left.\frac{d\xi_2}{dt}(t)\right|_{u=s=0},
\]
最后一个等式是因为 $\xi_2$ 是沿常环路 $x$ 的变分。综合可得
\[
 d^2\mathcal A_0([x,\widetilde x])(\xi_1,\xi_2)
 =-\int_0^1\omega\left(\frac{d\xi_2}{dt}(t),\xi_1(t)\right)dt.
\]
特别地，$d^2\mathcal A_0([x,\widetilde x])$ 的核由沿常道路 $\widetilde x$ 的向量场
\[
 \left\{\xi_2\in C^\infty(S^1,T_xM)\ \middle|\
 \frac{d\xi_2}{dt}=0\right\}
\]
组成。通过嵌入 $M\hookrightarrow C^\infty(S^1,M)$，可以把 $T_xM$ 与
$C^\infty(S^1,T_xM)$ 中的像 $T_xM$ 认同。因此
\[
 \ker d^2\mathcal A_0([x,\widetilde x])\cong T_xM,
\]
也就是说，该核恰好是 $\mathcal A_0$ 的临界流形
$\operatorname{Crit}\mathcal A_0$ 在任意点
$[x,\widetilde x]\in\operatorname{Crit}\mathcal A_0$ 处的切空间。对任意
$A\in\Gamma$，其他临界点 $[x,\widetilde x\#A]$ 也同样成立。
\end{proof}

按此命题，$\mathcal A_0$ 可视为\personBott{}--\personMorse{}函数 (Bo54)。我们将看到，加入通有的
含时哈密顿函数 $H:[0,1]\times M\to\mathbb R$ 后，可以消去这些临界点中的绝大
多数，只留下 $H$ 的有限个\emph{非退化} 1-周期哈密顿轨道。\personArnold{}的著名猜想 (Ar65)
断言，相伴哈密顿流至少有 $SB(M)$（即 $H^*(M)$ 的秩）个（可缩）1-周期轨道。

\section{扰动作用量泛函及其作用谱}
\label{sec:perturbed-action-functionals-spectrum}

给定 1-周期哈密顿量 $H:(\mathbb R/\mathbb Z)\times M\to\mathbb R$ 后，考虑扰动泛函
$\mathcal A_H:\widetilde{\mathcal L}(M)\to\mathbb R$：
\begin{equation}
 \mathcal A_H([\gamma,w])
 =\mathcal A_0([\gamma,w])-\int H(t,\gamma(t))\,dt
 =-\int w^*\omega-\int H(t,\gamma(t))\,dt.
 \label{eq:18-perturbed-action-functional}
\end{equation}
除非另作说明，我们总取 1-周期\emph{均值归一化哈密顿量}，即满足
\begin{equation}
 \int_M H_t\,\omega^n=0\qquad\text{对所有 }t\in S^1.
 \label{eq:18-mean-normalized-hamiltonian}
\end{equation}
记 $\mathcal H_0(M)=C_0^\infty(S^1\times M)$ 为这类均值归一化哈密顿量的集合。

\begin{lemma}
\label{lem:18-critical-points-perturbed-action}
$\mathcal A_H$ 的临界点集为
\[
 \operatorname{Crit}(\mathcal A_H)
 =\{[z,w]\mid z\in\operatorname{Per}(H),\ \partial w=z\},
\]
而 $\Gamma$ 在 $\widetilde{\mathcal L}_0(M)$ 上的作用典范地限制到该集合上。
\end{lemma}

\begin{definition}
\label{def:18-action-spectrum}
\stfhIndex{I027}{Zuoyongpu}{作用谱（action spectrum）}
定义 $H$ 的\emph{\termActionSpectrum{}}为
\[
 \operatorname{Spec}(H):=\{\mathcal A_H([z,w])\in\mathbb R\mid
 [z,w]\in\widetilde{\mathcal L}_0(M),\ z\in\operatorname{Per}(H)\},
\]
即 $\mathcal A_H:\widetilde{\mathcal L}(M)\to\mathbb R$ 的临界值集合。对每个给定的
$z\in\operatorname{Per}(H)$，记
\[
 \operatorname{Spec}(H;z)=\{\mathcal A_H([z,w])\in\mathbb R\mid
 [z,w]\in\pi^{-1}(z)\}.
\]
\end{definition}

\stfhIndex{I009}{00-gammaomega-torsor}{$\Gamma_\omega$-挠子（$\Gamma_\omega$-torsor）}
注意，$\operatorname{Spec}(H;z)$ 是一个 $\Gamma_\omega$-\termTorsor{}
\footnote{译者注：原书称固定 $z$ 的 $\operatorname{Spec}(H;z)$ 是“位于
$\operatorname{Per}(H)$ 上的 $\Gamma_\omega$ 主丛”，但固定纤维不能同时是整个基底
上的丛；译文按随后不交并公式表述为位于 $z$ 上的挠子。}，位于
$z\in\operatorname{Per}(H)$ 之上。因此
\[
 \operatorname{Spec}(H)=\bigsqcup_{z\in\operatorname{Per}(H)}
 \operatorname{Spec}(H;z).
\]
由于 $\Gamma_\omega=\omega(\pi_2(M))$ 是 $\mathbb R$ 的可数子群，它要么离散，要么
在 $\mathbb R$ 中稠密。下面的结论是研究改变哈密顿量 $H$ 时
$\mathcal A_H$ 的临界值行为的重要工具。这里的证明取自 (Oh02)。

\begin{lemma}
\label{lem:18-action-spectrum-measure-zero}
设 $H$ 是任意周期哈密顿量。则对任意 $H$，$\operatorname{Spec}(H)$ 都是
$\mathbb R$ 的零测集。
\end{lemma}

\begin{proof}
首先，对每个给定的周期轨道 $z$，$\operatorname{Spec}(H;z)\subset\mathbb R$ 是可数
集，因为 $\Gamma_\omega$ 可数。在每个 $z\in\operatorname{Per}(H)$ 的一个管状邻域
中考虑\personPoincare{}回归映射。更确切地说，取 $z(0)$ 的一个小邻域 $V\subset M$，将
$V$ 与以 $z(0)$ 为球心的 $2n$ 维球 $B^{2n}(\delta)$ 认同。再取球邻域
$V'=B^{2n}(\delta')$，满足 $V\Subset V'$，使（第一）\personPoincare{}回归映射
\[
 R_z:V\longrightarrow V',\qquad p\longmapsto\phi_H^1(p)
\]
有定义。现在如下定义从 $V$ 到 $S^1\cong\mathbb R/\mathbb Z$ 到 $M$ 的分片光滑映射
空间的连续映射：对每个 $p\in V$，先沿 $X_H$ 的流前进，再在 $V'$ 与
$B^{2n}(\delta')$ 的认同下沿直线从 $R_z(p)$ 返回 $p$。最后缩放环路的定义域，将其
重新参数化为 $[0,1]$。

记 $z_p$ 为上述构造中与 $p\in V$ 对应的环路，$\mathcal V_z\subset\mathcal L_0(M)$
为赋值 $p\mapsto z_p$ 的像。显然 $z_p$ 与 $z$ 同伦，因此任何以 $z$ 为边界的给定
圆盘 $w$ 都可自然延拓成以 $z_p$ 为边界的圆盘，记为 $w_p$。容易验证函数
\[
 h:\pi^{-1}(\mathcal V_z)\longrightarrow\mathbb R,
 \qquad h([z_p,w_p]):=\mathcal A_H([z_p,w_p])
\]
在 $\pi^{-1}(\mathcal V_z)$ 上光滑，并且它的临界值包含
$\operatorname{Spec}(H;z)$ 附近的 $\mathcal A_H$ 的临界值。这可通过写出
$\mathcal A_H([z_p,w_p])$ 的显式表达式并作简单的局部计算来证明。注意
$\pi^{-1}(\mathcal V_z)$ 是有限维（事实上为 $2n$ 维）流形，由\personSard{}定理，临界值集
在 $\mathbb R$ 中是零测集。所有 $z\in\operatorname{Per}(H)$ 的初始点 $z(0)$ 组成
紧集，因此有限多个这样的管状邻域连同其补集覆盖 $M$。故
$\operatorname{Spec}(H)\subset\mathbb R$ 包含在有限个零测集的并中，从而本身也是
零测集。
\end{proof}

\begin{exercise}
\label{ex:18-smoothness-local-action-function}
证明上面证明中定义的函数 $h$ 确实光滑。
\end{exercise}

注意，当 $H=0$ 时，
\[
 \operatorname{Spec}(H)=\Gamma_\omega.
\]

\begin{definition}
\label{def:18-homotopic-hamiltonians}
若 $\phi_H^1=\phi_F^1$，并且相伴的哈密顿路径
$\phi_H,\phi_F\in\mathcal P(\operatorname{Ham}(M,\omega),\mathrm{id})$ 保端点道路同伦，
则称两个哈密顿量 $H$ 与 $F$ \emph{同伦}。此时记 $H\sim F$，并以
$\widetilde{\operatorname{Ham}}(M,\omega)$ 表示等价类的集合。
\footnote{译者注：原书在定义 18.3.5 中把第二条相伴路径印作 $\phi_K$，但定义的两个
哈密顿量是 $H$ 与 $F$；译文按上下文改为 $\phi_F$。}
\end{definition}

下面的命题是研究哈密顿量形变下作用量泛函临界值的另一重要工具。它在非球面情形
由 (Schw00)、(Po01) 证明，在一般情形由 (Oh05a) 证明。

\begin{proposition}
\label{prop:18-homotopic-hamiltonians-same-spectrum}
设 $F,G$ 均为均值归一化哈密顿量。若 $F\sim G$，则作为 $\mathbb R$ 的子集，
\[
 \operatorname{Spec}F=\operatorname{Spec}G.
\]
\end{proposition}

证明前先作些准备。设 $\phi$ 是给定的哈密顿微分同胚，
$\{F^s\}_{s\in[0,1]}$ 是一条哈密顿量道路，满足
$F^0=G$、$F^1=F$，且对所有 $s\in[0,1]$ 都有 $\phi_{F^s}^1=\phi$。考虑哈密顿
微分同胚的二参数族
\[
 \phi(s,t)=\phi_{F^s}^t\circ(\phi_G^t)^{-1}.
\]
记 $F(s,t,x)=F^s(t,x)$，并令 $K=K(s,t,x)$ 是 $s$ 方向上的哈密顿函数，即与向量场
\[
 \frac{\partial\phi_{F^s}^t}{\partial s}\circ(\phi_{F^s}^t)^{-1}
\]
相伴的哈密顿量。若要求 $K(s,t,\cdot)$ 归一化，则 $K$ 唯一，而这种归一化总能实现。
此外，由 (2.4.26)，
\[
 \frac{\partial X_F}{\partial s}-\frac{\partial X_K}{\partial t}
 +[X_K,X_F]=0,
\]
继而由 (2.2.9) 和引理 2.3.4，
\[
 \frac{\partial F}{\partial s}-\frac{\partial K}{\partial t}-\{K,F\}=c(s,t),
\]
其中 $c$ 只依赖于 $(s,t)$。由于 $F$ 与 $K$ 都已归一化，而由\personLiouville{}引理，\personPoisson{}括号 $\{K,F\}$ 自动归一化，所以 $c\equiv0$。换言之，$F,K$ 满足
\begin{equation}
 \frac{\partial F}{\partial s}-\frac{\partial K}{\partial t}-\{K,F\}=0.
 \label{eq:18-hamiltonian-two-parameter-identity}
\end{equation}
值得把下面的等价恒等式单列出来。

\begin{lemma}
\label{lem:18-equivalent-two-parameter-identity}
等式~\eqref{eq:18-hamiltonian-two-parameter-identity} 等价于
\begin{equation}
 \frac{\partial F}{\partial s}(s,t,\phi_{F^s}^t(p))
 =\frac\partial{\partial t}\bigl(K(s,t,\phi_{F^s}^t(p))\bigr).
 \label{eq:18-hamiltonian-total-derivative}
\end{equation}
\end{lemma}

\begin{proof}
假设~\eqref{eq:18-hamiltonian-two-parameter-identity} 成立，则直接计算
\begin{align*}
 \frac{\partial F}{\partial s}(s,t,\phi_{F^s}^t(p))
 &=\frac{\partial K}{\partial t}(s,t,\phi_{F^s}^t(p))
   +\{K,F\}(s,t,\phi_{F^s}^t(p))\\
 &=\frac{\partial K}{\partial t}(s,t,\phi_{F^s}^t(p))
   +dK(X_F)(\phi_{F^s}^t(p))\\
 &=\frac\partial{\partial t}\bigl(K(s,t,\phi_{F^s}^t(p))\bigr),
\end{align*}
这就证明了~\eqref{eq:18-hamiltonian-total-derivative}。反过来，只需从右端的最后一行
倒着读回第一行即可。证明完毕。
\end{proof}

现在可以证明命题~\ref{prop:18-homotopic-hamiltonians-same-spectrum}。

\begin{proof}[命题~\ref{prop:18-homotopic-hamiltonians-same-spectrum} 的证明]
设 $\phi$ 是给定的哈密顿微分同胚，取满足上述条件的道路
$\{F^s\}_{s\in[0,1]}$，并使用二参数族
$\phi(s,t)=\phi_{F^s}^t\circ(\phi_G^t)^{-1}$ 及哈密顿量 $F,K$。注意，对每个
$s\in[0,1]$，$t\mapsto\phi(s,t)$ 都定义一条满足 $\phi(s,0)\equiv\mathrm{id}$ 的
可缩哈密顿环路，这是关于 $\{F^s\}$ 的假设所致。记此环路为 $h^s$，记其满足
$\widetilde h^s(0)=\mathrm{id}$ 的提升为 $\widetilde h^s$，它位于
$\widetilde{\operatorname{Ham}}(M,\omega)$ 中。环路 $\widetilde h^s$ 给出一一对应
\footnote{译者注：原书把下式的目标写作 $\operatorname{Crit}\mathcal A_F$；但
$h^s$ 把 $G$ 的轨道送到 $F^s$ 的轨道，且紧接着定义的 $\chi(s)$ 也使用
$\mathcal A_{F^s}$，故译文补上上标 $s$。}
\[
 (\widetilde h^s)_*: \operatorname{Crit}\mathcal A_G
 \longrightarrow\operatorname{Crit}\mathcal A_{F^s},
\]
其公式为
\[
 \widetilde h^s\cdot[z,w]=[h^s\cdot z,\widetilde h^s\cdot w].
\]
因此只需证明函数
\[
 \chi:[0,1]\to\mathbb R,\qquad
 \chi(s)=\mathcal A_{F^s}(\widetilde h^s\cdot[z,w])
\]
为常数。

由于对所有 $s\in[0,1]$ 都有
$\phi_{F^s}^1\equiv\phi_G^1$、$\phi_{F^s}^0\equiv\mathrm{id}$，故对每个 $s$，
$K(s,1,x)$ 和 $K(s,0,x)$ 都是常数；由归一化条件进一步得到
\begin{equation}
 K(s,1,\cdot)=K(s,0,\cdot)\equiv0.
 \label{eq:18-k-endpoint-normalization}
\end{equation}
对 $\chi(s)$ 求导：
\[
 \chi'(s)=d\mathcal A_{F^s}(\widetilde h^s\cdot[z,w])
 \left(\frac\partial{\partial s}(\widetilde h^s\cdot[z,w])\right)
 -\int_0^1\frac{\partial F}{\partial s}
 (s,t,(h^s\cdot z)(t))\,dt.
\]
因为 $\widetilde h^s\cdot[z,w]$ 是 $\mathcal A_{F^s}$ 的临界点，第一项为零。对第二项，
注意 $z$ 是 $\dot x=X_G(t,x)$ 的周期解，故存在 $\phi_G^1$ 的不动点 $p\in M$，使
$z(t)=\phi_G^t(p)$。于是
\[
 (h^s\cdot z)(t)=h_t^s(z(t))
 =(\phi_{F^s}^t\circ(\phi_G^t)^{-1})\circ\phi_G^t(p)
 =\phi_{F^s}^t(p).
\]
从而
\[
 \int_0^1\frac{\partial F}{\partial s}(s,t,(h^s\cdot z)(t))\,dt
 =\int_0^1\frac{\partial F}{\partial s}(s,t,\phi_{F^s}^t(p))\,dt.
\]
代入~\eqref{eq:18-hamiltonian-total-derivative}，在 $[0,1]$ 上积分全导数，并使用
\eqref{eq:18-k-endpoint-normalization}，得到
\[
 \int_0^1\frac{\partial F}{\partial s}(s,t,(h^s\cdot z)(t))\,dt
 =K(s,1,\phi_{F^s}^1(p))-K(s,0,\phi_{F^s}^0(p))=0.
\]
因此 $\chi$ 为常数，命题得证。
\footnote{译者注：原书本式末端两处写成未定义的 $f_1^s(p)$、$f_0^s(p)$；译文按本证明
始终使用的哈密顿流补为 $\phi_{F^s}^1(p)$、$\phi_{F^s}^0(p)$。}
\end{proof}

这个命题说明，归一化哈密顿量 $H$ 的作用谱 $\operatorname{Spec}H$ 只依赖于它的
同伦类。\footnote{译者注：原书此处称“这个引理”，但给出谱不变性结论的是命题
18.3.6；引理 18.3.7 只是其证明中使用的二参数恒等式，译文据此改为“这个命题”。}

\begin{definition}
\label{def:18-spectrum-hamiltonian-path-class}
设 $h\in\widetilde{\operatorname{Ham}}(M,\omega)$。取满足
$h=[\phi,F]$ 的某个（因而任意一个）归一化哈密顿量 $F$，定义 $h$ 的谱为
\[
 \operatorname{Spec}(h):=\operatorname{Spec}F.
\]
\end{definition}

显然，无论 $H$ 是否归一化，作用量泛函 $\mathcal A_H$ 都有定义，其作用谱亦然。记
相应的归一化哈密顿量为 $\underline H$，即
\begin{equation}
 \underline H(t,x)=H(t,x)-\frac1{\operatorname{vol}_\omega(M)}
 \int_M H_t(x)\,d\mu_\omega,
 \label{eq:18-hamiltonian-normalization}
\end{equation}
则\footnote{译者注：原书式 (18.3.21) 的右端写成
$\operatorname{Spec}(H)-\operatorname{Cal}(H)$。但由式 (18.3.15) 的约定
$\mathcal A_H=\mathcal A_0-\int H\,dt$ 及式 (18.3.20) 的
$\underline H=H-\operatorname{Cal}$，直接代入得到下式中的加号；译文据此改正。}
\begin{equation}
 \operatorname{Spec}(\underline H)
 =\operatorname{Spec}(H)+\operatorname{Cal}(H),
 \label{eq:18-spectrum-normalization-shift}
\end{equation}
其中
\[
 \operatorname{Cal}(H):=\frac1{\operatorname{vol}_\omega(M)}
 \int_0^1\int_M H_t(x)\,d\mu_\omega\,dt.
\]

\section{\texorpdfstring{$[z,w]$ 的\personConley{}--\personZehnder{}指标}{[z,w] 的康利--曾德尔指标}}
\label{sec:conley-zehnder-index-capped-orbit}
\stfhSubIndexRangeStart{I065}{Kangli-Zengdeerzhibiao}{康利--曾德尔指标（Conley--Zehnder index）}{Zzwdezhibiao}{$[z,w]$ 的指标（of $[z,w]$）}

首先注意，辛群 $\operatorname{Sp}(S,\Omega)\cong\operatorname{Sp}(\mathbb R^{2n})$
有极大紧子群 $U(n)$，$\operatorname{Sp}(S,\Omega)$ 形变收缩到该子群
\footnote{译者注：原书在本段形变收缩句及下一页共轭作用句中，把同一辛向量空间的
形式从 $\Omega$ 改写成 $\omega$；译文按本节开头的定义统一为
$\operatorname{Sp}(S,\Omega)$。}。因此行列式映射
$\det:U(n)\to S^1$ 诱导自然同态
\[
 \mu:\pi_1(\operatorname{Sp}(S,\Omega),\mathrm{id})\longrightarrow\mathbb Z.
\]
接下来进一步展开第 2.2.2 节关于 $\operatorname{Sp}(S,\Omega)$ 结构的讨论。使用
(SZ92) 的记号，令
\[
 \operatorname{Sp}(S,\Omega)^*=
 \{A\in\operatorname{Sp}(S,\Omega)\mid\det(I-A)\neq0\}.
\]
注意，当 $p\in\operatorname{Fix}\phi_H^1$ 非退化时，
$d\phi_H^1(p)\in\operatorname{Sp}(T_pM,\omega_p)$。下面由 (CZ84) 证明的引理对后续
讨论十分重要；其证明留给原文或 (SZ92)。

\begin{lemma}
\label{lem:18-symplectic-group-star-components}
$\operatorname{Sp}(S,\Omega)^*$ 有两个连通分支
\[
 \operatorname{Sp}^{\pm}(S,\Omega)=
 \{A\in\operatorname{Sp}(S,\Omega)^*\mid
   \pm\det(I-A)>0\}.
\]
此外，$\operatorname{Sp}(S,\Omega)^*$ 中的每个环路在
$\operatorname{Sp}(S,\Omega)$ 中可缩。
\end{lemma}

因此，$\operatorname{Sp}(S,\Omega)^*$ 和
$\operatorname{Sp}(S,\Omega)^\pm$ 在 $\operatorname{Sp}(S,\Omega)$ 的共轭作用下不变。
考虑道路集
\begin{equation}
 \operatorname{SP}^*(1)=\{\alpha:[0,1]\to\operatorname{Sp}(S,\Omega)
 \mid\alpha(0)=\mathrm{id},\ \alpha(1)\in\operatorname{Sp}^*(S,\Omega)\}.
 \label{eq:18-sp-star-path-space}
\end{equation}
不失一般性，设 $\alpha(1)\in\operatorname{Sp}^+(2n;\mathbb R)$；
$\operatorname{Sp}^-(2n;\mathbb R)$ 的情形相同。由引理~\ref{lem:18-symplectic-group-star-components}，
可选取同伦意义下唯一的道路
$\alpha_+^{\mathrm{can}}:[1,2]\to\operatorname{Sp}(2n;\mathbb R)$，使
\begin{equation}
 \begin{gathered}
  \alpha_+^{\mathrm{can}}(1)=\alpha(1),\qquad
  \alpha_+^{\mathrm{can}}(2)=W^+,\\
  \alpha_+^{\mathrm{can}}(t)\in\operatorname{Sp}^+(2n;\mathbb R)
  \quad\text{对所有 }t\in[1,2],
 \end{gathered}
 \label{eq:18-canonical-extension-path}
\end{equation}
其中 $W^+=-\mathrm{id}$。
\footnote{译者注：原书把 $\alpha_+^{\mathrm{can}}$ 的定义域印作 $[0,1]$，但紧接着
指定其在 $t=1,2$ 的端值，并要求 $t\in[1,2]$；译文据此改为 $[1,2]$。}
对 $\alpha(1)\in\operatorname{Sp}^-(2n)$ 的道路，考虑矩阵
\[
 W^-=\operatorname{diag}(2,-1,\ldots,-1,1/2,\ldots,-1)
 \in\operatorname{Sp}^-(2n),
\]
并类似地定义 $\alpha_-^{\mathrm{can}}$。于是 (CZ84) 构造了下面的连续映射
$\rho:\operatorname{Sp}(V,\omega)\to S^1$；证明亦见 (SZ92) 的定理 3.1。

\begin{theorem}
\label{thm:18-salamon-zehnder-rho-map}
设 $(V,\omega)$ 是辛向量空间。则存在唯一的一族连续映射
\[
 \rho:\operatorname{Sp}(V,\omega)\longrightarrow S^1
\]
（每个 $(V,\omega)$ 一个），满足：
\begin{enumerate}[label=(\arabic*)]
\item \stfhPropertyLabel{（自然性）} 若 $T:(V_1,\omega_1)\to(V_2,\omega_2)$ 是辛同构，则对所有
$A\in\operatorname{Sp}(V_1,\omega_1)$，
$\rho(TAT^{-1})=\rho(A)$。
\item \stfhPropertyLabel{（乘积）} 若 $(V,\omega)=(V_1\times V_2,\omega_1\times\omega_2)$，
且 $A=A_1\oplus A_2\in\operatorname{Sp}(V,\omega)$，则
$\rho(A)=\rho(A_1)\rho(A_2)$。
\item \stfhPropertyLabel{（行列式）} 若
$A\in\operatorname{Sp}(2n,\mathbb R)\cap O(2n)=U_{\mathbb R}(n)$，则
$\rho(A)=\det(X+iY)$，其中
\[
 A=\begin{pmatrix}X&-Y\\Y&X\end{pmatrix},\qquad X+iY\in U(n).
\]
\item \stfhPropertyLabel{（归一化）} 若 $A$ 在单位圆上没有特征值，则 $\rho(A)=\pm1$。
\end{enumerate}
\end{theorem}

把 $\alpha$ 与 $\alpha_+^{\mathrm{can}}$ 连接，并考虑复合
\[
 \rho\circ(\alpha\#\alpha_+^{\mathrm{can}}):[0,2]\longrightarrow S^1,
\]
其端点值为 $\pm1$；因而将它平方后得到一个环路，其环绕数由映射
$\det^2:U(n)\to S^1$ 诱导。

\begin{definition}
\label{def:18-conley-zehnder-index-path}
设 $\alpha\in\operatorname{SP}^*(1)$。定义 $\alpha$ 的\personConley{}--\personZehnder{}指标为
\footnote{译者注：原书定义 18.4.3 印作
$\deg\det^2(\rho\circ(\alpha\#\alpha_+^{\mathrm{can}}))$；但 $\rho$ 的值已在
$S^1$ 中，不能再施加定义在 $U(n)$ 上的 $\det$。结合前句及 (SZ92) 的定义，译文把
它写成复合映射的平方之次数。}
\[
 \mu_{\mathrm{CZ}}(\alpha)
 =\deg\bigl((\rho\circ(\alpha\#\alpha_+^{\mathrm{can}}))^2\bigr).
\]
\end{definition}

定义~\ref{def:18-conley-zehnder-index-path} 的 $\mu_{\mathrm{CZ}}$ 与 (SZ92) 的定义
一致。另一方面，对给定\personLie{}群 $G$，记
\[
 \Omega(G)=\{g:S^1\to G\mid g(0)=\mathrm{id}\}
\]
为 $G$ 的基环路群。于是有明显的指标映射
$\mu:\Omega(\operatorname{Sp}(2n,\mathbb R))\to\mathbb Z$，它由形变收缩
$\operatorname{Sp}(2n,\mathbb R)\simeq U(n)$ 及复合映射的次数定义；具体地，对
$g\in\Omega(U(n))$，该整数由
\[
 g\longmapsto\frac1{2\pi}\int_{S^1}(\det\circ g)^*d\theta
\]
给出。
\footnote{译者注：原书在这里把相应同伦群印作
$\pi_1(\Omega(\operatorname{Sp}(2n,\mathbb R)),\mathrm{id})$ 与
$\pi_1(\Omega(U(n)),\mathrm{id})$，这会多取一层环路空间，且与紧接着对
$g:S^1\to\operatorname{Sp}(2n,\mathbb R)$ 的积分定义不相容；译文按后者陈述。}
\personConley{}--\personZehnder{}指标 $\mu_{\mathrm{CZ}}$ 可由下列性质刻画。

\begin{proposition}[{(FH93) 中的命题 5}]
\label{prop:18-conley-zehnder-characterization}
存在唯一映射
\[
 \mu_{\mathrm{CZ}}:\operatorname{SP}^*(1)\longrightarrow\mathbb Z
\]
满足以下准则：
\begin{enumerate}[label=(\arabic*)]
\item \stfhPropertyLabel{（归一化）}
\[
 \mu_{\mathrm{CZ}}\!\left(\tau\mapsto
 \{e^{-\tau}+ie^\tau\}\oplus\{e^{\pi i\tau}\}\oplus\cdots\oplus
 \{e^{\pi i\tau}\}\right)=n-1,
 \qquad
 \mu_{\mathrm{CZ}}(\tau\mapsto e^{\pi i\tau}\mathrm{Id})=n.
\]
\item \stfhPropertyLabel{（基环路作用）} 在基环路群
$\Omega(\operatorname{Sp}(2n,\mathbb R),\mathrm{id})$ 对
$\operatorname{SP}^*(1)$ 的作用下，对每个环路
$g:S^1\to\operatorname{Sp}(2n,\mathbb R)$，
\[
 \mu_{\mathrm{CZ}}(g\alpha)=2\mu(g)+\mu_{\mathrm{CZ}}(\alpha).
\]
\end{enumerate}
\end{proposition}

取 $[z,w]\in\operatorname{Crit}\mathcal A_H$，并选取\termTrivialization{}
\[
 \Psi:w^*TM\longrightarrow D^2\times\mathbb R^{2n}.
\]
回忆 $z(t)=\phi_H^t(p)$，其中 $p\in\operatorname{Fix}\phi_H^1$，于是可写
\[
 \Psi\circ d\phi_H^t\circ\Psi^{-1}(t,v)=(t,\alpha_\Psi(t)v).
\]
按定义，
\[
 \alpha_\Psi(0)=\mathrm{id},\qquad
 \alpha_\Psi(1)\in\operatorname{Sp}^*(2n;\mathbb R).
\]

\begin{definition}
\label{def:18-conley-zehnder-capped-orbit}
定义 $[z,w]$ 的\personConley{}--\personZehnder{}指标 $\mu_H([z,w])$ 为
\[
 \mu_H([z,w])=\mu_{\mathrm{CZ}}(\alpha_\Psi).
\]
\end{definition}

下面陈述“粘上一个球面”作用
$[z,w]\mapsto[z,w\#A]$ 下\personConley{}--\personZehnder{}指标与\termFirstChernNumber{}
$c_1(A)$ 之间的重要恒等式。其指标公式的证明沿用 (Oh06a) 的论述。

\begin{theorem}
\label{thm:18-conley-zehnder-capping-change}
设 $z:S^1=\mathbb R/\mathbb Z\to M$ 是
$\dot x=X_H(x)$ 的给定 1-周期解，$w,w'$ 是两个给定的填充圆盘。则
\begin{equation}
 \mu_H([z,w'])=\mu_H([z,w])-2c_1([w'\#\overline w]).
 \label{eq:18-index-two-cappings}
\end{equation}
特别地，
\begin{equation}
 \mu_H([z,w\#A])=\mu_H([z,w])-2c_1(A).
 \label{eq:18-index-sphere-gluing}
\end{equation}
\end{theorem}

\begin{remark}
\label{rem:18-conley-zehnder-sign-convention}
要强调的是，在我们的约定下，公式中第一\personChern{}数前的符号是“$-$”。它与 (HS95)
公式的符号不同，源于 $\mathbb C^n$ 上典范辛形式的约定不同：把
$\mathbb R^{2n}\cong T^*\mathbb R^n$，并把相应典范坐标记为
$(q_1,\ldots,q_n,p_1,\ldots,p_n)$ 时，我们约定的典范辛形式为
\begin{equation}
 \omega_0=\sum_jdq_j\wedge dp_j,
 \label{eq:18-canonical-symplectic-sign}
\end{equation}
而 (HS95)、(SZ92) 或 (Po01) 采用
\[
 \omega_0'=-\omega_0=\sum_jdp_j\wedge dq_j.
\]
\end{remark}

\begin{proof}
把指标公式的证明分成若干步。

\begin{enumerate}[label=\arabic*.]
\item \stfhPropertyLabel{（典范辛形式）} 在坐标 $z_j=q_j+ip_j$ 下，我们对
$T^*\mathbb R^n=\mathbb R^{2n}\cong\mathbb C^n$ 的典范辛形式采用
\eqref{eq:18-canonical-symplectic-sign} 的约定。

\item \stfhPropertyLabel{（典范复结构）} 设 $J_0$ 是 $\mathbb R^{2n}\cong\mathbb C^n$ 上由
复数 $i$ 的乘法给出的标准复结构，其中 $z_j=q_j+ip_j$。在我们对 $\mathbb C^n$ 上
典范辛形式 $\omega_0$ 的约定下，相伴的\personHermitian{}结构
\[
 \langle\cdot,\cdot\rangle:\mathbb C^n\times\mathbb C^n\longrightarrow\mathbb C
\]
\emph{对第一个变元复线性、对第二个变元反线性}。换言之，\personHermitian{}内积为
\begin{equation}
 \langle u,v\rangle=g(u,v)-i\omega_0(u,v),
 \label{eq:18-hermitian-our-convention}
\end{equation}
其中 $g$ 是 $\mathbb R^{2n}$ 上的标准欧氏内积。要注意，这个\personHermitian{}结构是
(HS95)、(SZ92)、(Po01) 所用结构的共轭；后者对应于
\begin{equation}
 \langle u,v\rangle=g(u,v)+i\omega_0(u,v).
 \label{eq:18-hermitian-other-convention}
\end{equation}
见 (SZ92) 引理 5.1 之前的注。等价地，后一个\personHermitian{}结构与殆\personKahler{}结构
\[
 (g,\omega_0',J_0')
\]
相伴，其中 $J_0'$ 是 $J_0$ 的共轭殆复结构。\emph{$\mathbb C^n$ 上复结构的这种改变，会
改变一般复向量丛 $E$ 的第一\personChern{}数的符号}；回忆复向量丛\personChern{}类的一般公式
（见 (MSt74) 引理 14.9）：
\[
 c_k(\overline E)=(-1)^kc_k(E).
\]

\item \stfhPropertyLabel{（$\operatorname{SP}^*(1)$ 上的\personConley{}--\personZehnder{}指标）}
注意，两种约定下 $\operatorname{Sp}(2n,\mathbb R)$ 的定义相同。

\item 给定两个映射 $w,w':D^2\to M$，满足
$w|_{\partial D^2}=w'|_{\partial D^2}$，如下定义粘合映射
$u=w\#\overline{w'}:S^2\to M$：
\[
 u(z)=\begin{cases}
  w(z),&z\in D^+,\\
  w'(1/\overline z),&z\in D^-.
 \end{cases}
\]
这里 $D^+$ 是保持原定向的 $D^2$，$D^-$ 带相反定向。先验地这只给出连续映射，
但通过在边界附近把映射“压平”，可在不改变同伦类的前提下将其形变为光滑映射
\footnote{译者注：原书说要把两个映射在边界附近“压平”，但紧接着的公式只写了
$w$。要使 $w\#\overline{w'}$ 沿粘合圆周光滑，两个圆盘都须作同样处理；译文补出
$w'$ 的条件。}。
换言之，可假设充分小的 $\epsilon>0$ 使得
\[
 \begin{gathered}
 w(z)=w(z/|z|),\qquad w'(z)=w'(z/|z|)\\
 \text{当 }|z|\geq1-\epsilon.
 \end{gathered}
\]
以下总假设填充圆盘在这个意义下是平坦的。作此调整后，$u$ 定义 $S^2$ 上的光滑映射。

\item \stfhPropertyLabel{（标记条件）} 对给定的 $[z,w]$、$[z,w']$，其中周期轨道
$z(t)=\phi_H^t(z(0))$，对平凡化
\[
 \Phi_w:w^*TM\to D^2\times(\mathbb R^{2n},\omega_0),\qquad
 \Phi_{w'}:(w')^*TM\to D^2\times(\mathbb R^{2n},\omega_0),
\]
附加总能满足的\emph{标记}条件
\begin{equation}
 \Phi_w(z(0))=\Phi_{w'}(z(0))
 \label{eq:18-trivialization-marking-condition}
\end{equation}
（视为从 $T_{z(0)}M$ 到 $\mathbb R^{2n}$ 的映射）。
\footnote{译者注：原书把 $\Phi_w$ 与 $\Phi_{w'}$ 的定义域都印作 $w^*TM$；后者按
其下标和后续使用应为 $(w')^*TM$，译文据此补正。}
\emph{在此附加条件下}，可以写成
\begin{equation}
 \alpha_{[z,w']}(t)=S_{w'w}(t)\cdot\alpha_{[z,w]}(t),
 \label{eq:18-capping-trivialization-loop}
\end{equation}
其中 $S_{w'w}:S^1=\mathbb R/\mathbb Z\to\operatorname{Sp}(2n,\mathbb R)$ 是由
\eqref{eq:18-capping-trivialization-loop} 定义的\emph{环路}。它确实是环路，因为
\begin{equation}
 \alpha_{[z,w']}(0)=\alpha_{[z,w]}(0)=\mathrm{id},
 \label{eq:18-capping-path-start}
\end{equation}
\begin{equation}
 \alpha_{[z,w']}(1)=\alpha_{[z,w]}(1).
 \label{eq:18-capping-path-end}
\end{equation}
事实上，由~\eqref{eq:18-capping-trivialization-loop} 与
\eqref{eq:18-trivialization-marking-condition}，有恒等式
\begin{align}
 S_{w'w}(t)
 &=\bigl(\Phi_{w'}(z(t))\circ d\phi_H^t(z(0))
          \circ\Phi_{w'}(z(0))^{-1}\bigr)\nonumber\\
 &\quad\circ\bigl(\Phi_w(z(t))\circ d\phi_H^t(z(0))
          \circ\Phi_w(z(0))^{-1}\bigr)^{-1}\nonumber\\
 &=\Phi_{w'}(z(t))\circ d\phi_H^t(z(0))\nonumber\\
 &\quad\circ\Phi_{w'}(z(0))^{-1}\circ\Phi_w(z(0))\nonumber\\
 &\quad\circ(d\phi_H^t)^{-1}(z(0))\circ(\Phi_w(z(t)))^{-1}.
 \label{eq:18-transition-loop-expanded}
\end{align}
\stfhIndex{I123}{Biaojitiaojian}{标记条件（marking condition）}
标记条件~\eqref{eq:18-trivialization-marking-condition} 使
\eqref{eq:18-transition-loop-expanded} 中间的项相消，因而
\begin{equation}
 S_{w'w}(t)=\Phi_{w'}(z(t))\circ\Phi_w(z(t))^{-1}.
 \label{eq:18-transition-loop}
\end{equation}
由命题~\ref{prop:18-conley-zehnder-characterization} 及
\eqref{eq:18-transition-loop}，得到
\begin{equation}
 \operatorname{ind}_{\mathrm{CZ}}(\overline\alpha_{[z,w']})
 =2\operatorname{wind}(\widehat S_{w'w})
  +\operatorname{ind}_{\mathrm{CZ}}(\overline\alpha_{[z,w]}).
 \label{eq:18-index-transition-winding}
\end{equation}
这里 $\widehat S_{w'w}:S^1\to U(n)$ 是在
$\operatorname{Sp}(2n,\mathbb R)$ 内与 $S_{w'w}$ 同伦的酉矩阵环路。这样的同伦总存在，
并且在同伦意义下唯一，因为 $U(n)$ 是 $\operatorname{Sp}(2n,\mathbb R)$ 的形变收缩。

\stfhIndex{I079}{Diyichenlei}{第一陈类（first Chern class）}
\item \stfhPropertyLabel{（$c_1$ 的归一化）} 最后回忆 $S^2$ 上辛向量丛 $E$ 的第一\personChern{}类
$c_1$ 的定义。我们把\personChern{}类归一化，使
$S^2\cong\mathbb{CP}^1$ 的切丛具有第一\personChern{}数 $2$；这也与文献中的标准约定一致。
\emph{要指出的是，这种归一化与我们约定中式~\eqref{eq:18-hermitian-our-convention}
给出的 $\mathbb C^n$ 上\personHermitian{}结构相容}（见 (Mil65) 第 167 页）。

分解 $S^2=D^+\cup D^-$，并考虑辛平凡化
\[
 \Phi_+:E|_{D^+}\to D^2\times(\mathbb R^{2n},\omega_0),\qquad
 \Phi_-:E|_{D^-}\to D^2\times(\mathbb R^{2n},\omega_0).
\]
在我们约定的 $\mathbb C^n$ 上\personHermitian{}结构下，它们同伦于酉平凡化；在另一约定
下，则同伦于共轭酉平凡化。记转移矩阵环路为
\[
 \phi_{+-}:S^1\longrightarrow\operatorname{Sp}(2n,\mathbb R),
\]
它由方程
\[
 \Phi_+|_{S^1}\circ(\Phi_-|_{S^1})^{-1}(t,\xi)
 =(t,\phi_{+-}(t)\xi)
\]
确定，其中 $(t,\xi)\in E|_{S^1}$，且
$S^1=\partial D^+=\partial D^-$。按定义，在我们的约定下
\begin{equation}
 c_1(E)=\mu(\widehat\phi_{+-}).
 \label{eq:18-chern-transition-our-convention}
\end{equation}
等价地，
\begin{equation}
 c_1(E)=-\mu(\overline{\widehat\phi}_{+-}).
 \label{eq:18-chern-transition-conjugate}
\end{equation}

现在把它应用到 $u^*(TM)$，其中 $u=w\#\overline{w'}$，$\Phi_w,\Phi_{w'}$ 是
第 5 步给出的平凡化。由~\eqref{eq:18-transition-loop}，$S_{w'w}$ 是
$\Phi_w$ 与 $\Phi_{w'}$ 之间的转移矩阵环路。因此按定义，第一\personChern{}数
$c_1(u^*TM)$ 就是酉矩阵环路
\[
 \widehat S_{w'w}:t\longmapsto\widehat S_{w'w}(t),
 \qquad S^1\longrightarrow U(n)
\]
在我们约定的 $\mathbb C^n$ 上\personHermitian{}结构中的环绕数
$\mu(\widehat S_{w'w})$。容易验证，把这个环绕数用于 $S^2$ 的切丛时，其值确为
$2$，因此与我们采用的\personChern{}类约定一致。
\footnote{译者注：原书此处称 $\Phi_w,\Phi_{w'}$ 在第 4 步给出；两者实际在第 5 步
“标记条件”中引入，译文据此更正步骤号。}

\item \stfhPropertyLabel{（证明收尾）} 综合以上各步，特别是第 2、5、6 步，公式
\eqref{eq:18-index-transition-winding} 和
\eqref{eq:18-chern-transition-conjugate} 就化为所要证明的指标公式。
\end{enumerate}
\end{proof}

\section{哈密顿扰动\personCauchy{}--\personRiemann{}方程}
\label{sec:hamiltonian-perturbed-cauchy-riemann-equation}

为了对 $\mathcal A_H$ 作\personMorse{}理论，需要在
$\widetilde{\mathcal L}_0(M)$ 上给出度量。为此先在 $\mathcal L_0(M)$ 上定义度量，
再将其拉回 $\widetilde{\mathcal L}_0(M)$。注意，$M$ 上任何\personRiemann{}度量的
$S^1$-族 $\{g_t\}_{t\in S^1}$ 都在 $\mathcal L(M)$ 上诱导 $L^2$ 型度量
\begin{equation}
 \langle\!\langle\xi_1,\xi_2\rangle\!\rangle
 =\int_0^1g_t(\xi_1(t),\xi_2(t))\,dt
 \label{eq:18-loop-space-l2-metric}
\end{equation}
（$\xi_1,\xi_2\in T_\gamma\mathcal L(M)$）。在辛流形 $(M,\omega)$ 上，特别考虑由
\stfhIndex{I034}{Daikaileduliang}{殆凯勒度量（almost-Kähler metric）}
与辛形式 $\omega$ 相容的殆复结构一参数族 $J$ 所诱导的\emph{殆\personKahler{}度量}族
\[
 g_J=\omega(\,\cdot,J\,\cdot)
\]
及相伴范数 $|\cdot|_J$。给定 $J=\{J_t\}_{t\in S^1}$，它在 $\mathcal L(M)$ 上诱导
相伴的 $L^2$ 度量 $\langle\!\langle\cdot,\cdot\rangle\!\rangle_J$，也可写成
\begin{equation}
 \langle\!\langle\xi_1,\xi_2\rangle\!\rangle_J
 =\int_0^1\omega(\xi_1(t),J_t\xi_2(t))\,dt
 \label{eq:18-loop-space-compatible-l2-metric}
\end{equation}
（$\xi_1,\xi_2\in T_\gamma\mathcal L(M)$）。除非另作说明，从现在起 $J$ 总表示相容
殆复结构的 \emph{$S^1$-族}，并记
\stfhIndex{I010}{00-jomega}{$\mathfrak j_\omega$}
\[
 \mathfrak j_\omega:=C^\infty(S^1,\mathcal J_\omega).
\]
若以 $\operatorname{grad}_J\mathcal A_H$ 表示相伴的 $L^2$ 梯度向量场，则由
\eqref{eq:18-global-action-primitive} 与
\eqref{eq:18-loop-space-compatible-l2-metric} 直接计算得到
\begin{equation}
 \operatorname{grad}_J\mathcal A_H([\gamma,w])(t)
 =J_t\bigl(\dot\gamma(t)-X_H(t,\gamma(t))\bigr),
 \label{eq:18-action-l2-gradient}
\end{equation}
以后简写为 $J(\dot\gamma-X_H(\gamma))$。

\stfhSubIndexRangeEnd{I065}{Kangli-Zengdeerzhibiao}{康利--曾德尔指标（Conley--Zehnder index）}{Zzwdezhibiao}{$[z,w]$ 的指标（of $[z,w]$）}
\begin{exercise}
\label{ex:18-prove-action-gradient}
证明这个公式。
\end{exercise}

由此可知，该梯度可以投影到 $\mathcal L_0(M)$。因此，把道路
$u:\mathbb R\to\widetilde{\mathcal L}_0(M)$ 的\emph{负}梯度流方程投影到
$\mathcal L_0(M)$ 后，若把 $u$ 看作映射 $u:\mathbb R\times S^1\to M$，方程写成
\begin{equation}
 \frac{\partial u}{\partial\tau}
 +J\left(\frac{\partial u}{\partial t}-X_H(u)\right)=0.
 \label{eq:18-floer-perturbed-cr-equation}
\end{equation}
称它为与二元组 $(H,J)$ 相伴的\emph{\personFloer{} \termPerturbedCauchyRiemannEquation{}}，或简称
扰动\personCauchy{}--\personRiemann{}方程。

\personFloer{}理论在很大程度上依赖于研究形如
\eqref{eq:18-floer-perturbed-cr-equation} 的扰动\personCauchy{}--\personRiemann{}方程的\emph{有限能量}解
$u:\mathbb R\times S^1\to M$ 的模空间。相应的离壳能量泛函定义如下。

\begin{definition}[能量]
\label{def:18-floer-energy}
给定光滑映射 $u:\mathbb R\times S^1\to M$，定义 $u$ 的能量
$E_{(H,J)}(u)$ 为
\[
 E_{(H,J)}(u)=\frac12\int_{\mathbb R\times S^1}
 \left(\left|\frac{\partial u}{\partial\tau}\right|_{J_t}^2
 +\left|\frac{\partial u}{\partial t}-X_H(u)\right|_{J_t}^2\right)
 dt\,d\tau.
\]
\end{definition}

下面的命题是引理 12.2.4 的闭弦版本，证明与此前相同，故略去。

\begin{proposition}
\label{prop:18-finite-energy-asymptotic-subsequence}
设 $(M,\omega)$ 是任意辛流形，不必紧，$H:S^1\times M\to\mathbb R$ 是任意
哈密顿量。设 $u:\mathbb R\times S^1\to M$ 是
\eqref{eq:18-floer-perturbed-cr-equation} 的有限能量解，且其像在 $\operatorname{int}M$
中有紧闭包。则存在序列 $\tau_k\to\infty$（分别地，$\tau_k\to-\infty$），使环路
$z_k:=u(\tau_k)=u(\tau_k,\cdot)$ 在 $C^\infty$ 意义下收敛到\personHamilton{}方程
$\dot x=X_H(x)$ 的一个 1-周期解 $z:S^1\to M$。
\end{proposition}

对不必非退化的一般 $H$，记
\[
 \mathcal M(H,J)=\mathcal M(H,J;\omega)
\]
为~\eqref{eq:18-floer-perturbed-cr-equation} 的有限能量解的集合。

对~\eqref{eq:18-floer-perturbed-cr-equation} 的\emph{非自治}版本可以作类似讨论，现说明如下。
首先回忆
\[
 \mathcal H_0(M)=\{H:S^1\times M\to\mathbb R\mid H\text{ 均值归一化}\}.
\]
考虑 $\mathbb R$-族
\[
 \mathcal H_R:\mathbb R\to\mathcal H_0,\quad\tau\mapsto H(\tau),
 \qquad
 j_R:\mathbb R\to\mathfrak j_\omega,\quad\tau\mapsto J(\tau),
\]
它们渐近常值，即当 $|\tau|>R$（其中 $R$ 是充分大的常数）时，
\[
 H(\tau)=H^{\pm\infty},\qquad J(\tau)=j^{\pm\infty},
\]
其中 $H^{\pm\infty}\in\mathcal H_0$、$j^{\pm\infty}\in\mathfrak j_\omega$。
\footnote{译者注：源书印 $\mathbb J_\mathbb R\colon\mathbb R\to\mathcal J_\omega$、
$j^{\pm\infty},j\in\mathcal J_\omega$、$J_\alpha,J_\beta\in\mathcal J_\omega$；
译文统一作 $j_\mathbb R$、$j^{\pm\infty}$ 等，且陪域归 $\mathfrak j_\omega$——因
式 (18.5.42) 与 (18.5.44) 中的 $J(\tau)$ 须为 $S^1$ 族，并须与第 19 章
$(\mathcal H^\rho,\mathfrak j^\rho)$ 相衔接。}
与任意
这样的二元组相伴有~\eqref{eq:18-floer-perturbed-cr-equation} 的如下\emph{非自治}版本：
\begin{equation}
 \frac{\partial u}{\partial\tau}
 +J(\tau)\left(\frac{\partial u}{\partial t}-X_{H(\tau)}(u)\right)=0.
 \label{eq:18-nonautonomous-floer-equation}
\end{equation}
相伴能量泛函为
\[
 E_{(\mathcal H_R,j_R)}(u)=\frac12\int_{-\infty}^{\infty}\int_0^1
 \left(\left|\frac{\partial u}{\partial\tau}\right|_{J(\tau)}^2
 +\left|\frac{\partial u}{\partial t}-X_{H(\tau)}(u)\right|_{J(\tau)}^2\right)
 dt\,d\tau.
\]
记
\[
 \mathcal M(\mathcal H_R,j_R)=\mathcal M(\mathcal H_R,j_R;\omega)
\]
为~\eqref{eq:18-nonautonomous-floer-equation} 的有限能量解的集合。

下面是命题~\ref{prop:18-finite-energy-asymptotic-subsequence} 的类似结果。由于
$(\mathcal H,j)$ 渐近常值，其证明本质上相同。

\begin{proposition}
\label{prop:18-nonautonomous-asymptotic-subsequence}
设 $\mathcal H_R,j_R$ 如上，且 $u:\mathbb R\times S^1\to M$ 是
\eqref{eq:18-nonautonomous-floer-equation} 的有限能量解。则存在序列
$\tau_k\to\infty$，使环路 $z_k:=u(\tau_k)=u(\tau_k,\cdot)$ 在 $C^\infty$ 意义下
收敛到\personHamilton{}方程 $\dot x=X_{H^{+\infty}}(x)$ 的 1-周期解
$z:S^1\to M$。在 $-\infty$ 端也有相同结论，此时极限满足
$\dot x=X_{H^{-\infty}}(x)$。
\end{proposition}

这类渐近常值族的一种典型来源，是把 $[0,1]$ 上给定的光滑一参数族\emph{拉长}。

\begin{definition}
\label{def:18-boundary-flat-map}
设 $T$ 是任意拓扑空间。若连续映射 $f:[0,1]\to T$ 在边界
\stfhIndex{I044}{Bianjiepingtan}{边界平坦（boundary flat）}
$\partial[0,1]=\{0,1\}$ 的邻域内为常值，则称 $f$ \emph{边界平坦}。
\end{definition}

设 $\mathcal H:[0,1]\to\mathcal H_0(M)$ 是连接两个哈密顿量
$H_\alpha,H_\beta\in\mathcal H_0(M)$ 的同伦，
$j:[0,1]\to\mathfrak j_\omega$ 连接 $J_\alpha,J_\beta\in\mathfrak j_\omega$。记
\[
 \mathcal P(\mathfrak j_\omega):=C^\infty([0,1],\mathfrak j_\omega),
 \qquad
 \mathcal P(\mathcal H_0(M)):=C^\infty([0,1],\mathcal H_0(M)).
\]
定义一类函数 $\rho:\mathbb R\to[0,1]$：
\begin{equation}
 \rho(\tau)=\begin{cases}
 0,&\tau\leq0,\\
 1,&\tau\geq1,
 \end{cases}
 \qquad \rho'(\tau)\geq0.
 \label{eq:18-elongation-function}
\end{equation}
\stfhIndex{I072}{Lachanghanshu}{拉长函数（elongation function）}
称 $\rho$ 为\termElongationFunction{}。

每个这样的二元组 $(\mathcal H,j)$ 连同拉长函数 $\rho$，定义一对渐近常值
$\mathbb R$-族 $(\mathcal H^\rho,j^\rho)$：
\[
 \mathcal H_R=\mathcal H^\rho,\qquad j_R=j^\rho,
\]
其中 $\mathcal H^\rho=\{H^\rho(\tau)\}_{\tau\in\mathbb R}$ 是重新参数化的同伦
\[
 \tau\longmapsto H^\rho(\tau,t,x)=H(\rho(\tau),t,x).
\]
称 $\mathcal H^\rho$ 为 $\mathcal H$ 的 \emph{$\rho$-拉长}，或 $\mathcal H$ 的
\emph{$\rho$-拉长同伦}；对 $j$ 作相同定义。因此，这样的三元组 $(\mathcal H,j;\rho)$ 给出
\emph{非自治}方程
\begin{equation}
 \frac{\partial u}{\partial\tau}
 +J^\rho\left(\frac{\partial u}{\partial t}-X_{H^\rho}(u)\right)=0.
 \label{eq:18-elongated-floer-equation}
\end{equation}
记
\[
 \mathcal M(\mathcal H,j;\rho)
\]
为~\eqref{eq:18-elongated-floer-equation} 的有限能量解集合。

